% ===== §12 Exploitation in symbolic generation =====
\section{Exploitation in Symbolic Generation, and Generative Justice}
\label{sec:exploitation}

\noindent
The practice was specified as the cultivation of a field rather than the construction of a grammar: as the guidance of a generative dynamics rather than the modification of a generative structure (\xref{sec:field-not-grammar}). That specification was grounded formally, in the impossibility of constructing a structure one does not possess and which is itself in generation (\xref{sec:modelling}). It has, beyond its formal ground, an ethical content, and this section develops it. The thesis is that the possessive modification of a generative structure is a form of exploitation, that its avoidance is a form of generative justice, and that this brings the symbolic register of intimacy under the same criterion of justice the series applies to the economic. The section thereby connects the practice of \xref{sec:praxis} to the ethics of \xref{sec:keystone}, and supplies the second of the two forms in which the good may fail.

\subsection{From non-possession to the means of symbolic generation}
\label{sec:exp-means}

The practice does not construct the grammar; it guides the grammar's derivations. The ground given in \xref{sec:field-not-grammar} was that the grammar cannot be constructed, since to construct it would be to legislate the productions, and the meta-grammar that would have to be possessed in order to legislate them does not exist as a fixed possessable structure but is itself in generation and without a fixed ground (\xref{sec:modelling}). To construct the grammar is therefore to claim possession of a generativity that is no one's, to occupy the position of the barred Other as though it were the unbarred Other. This is the formal content of non-possession. Its ethical content is reached by a single further step.

The step is to identify the generative structure as a means of generation. Whoever possesses the generative structure of a relation, its grammar, the rules by which its value is generated, controls thereby the generation and the distribution of the relation's value. This is the symbolic instance of the general thesis of political economy, that the possession of the means of production confers the power to direct production and to appropriate its surplus \citep{marx1976}. The generative structure is the means of symbolic generation; its possession confers the power to direct what value the relation generates and to whom that value accrues.

\begin{claim}[The generative structure as a means of symbolic generation]
The grammar of a relation is a means of symbolic generation, in the sense in which the means of production is a means of material generation. To possess it, to occupy the position from which its productions are legislated, is to control the generation and the distribution of the relation's value. The possession of a generative structure is therefore not a neutral formal fact but the symbolic correlate of the possession of the means of production: it confers the power to direct generation and to appropriate the generated. The bridge from the possession of structure to the appropriation of value is the same bridge Marx draws in the material case: whoever controls the means by which value is generated controls the distribution of what is generated.
\end{claim}

\subsection{Possession as exploitation, guidance as generative justice}
\label{sec:exp-justice}

The criterion of generative justice is that value be retained by and returned to the generative network that produces it, rather than extracted from the network to an external possessor; injustice is the extraction of value from its generators \citep{eglash2016}. Applied to the symbolic register, the criterion yields the following. The possession of the generative structure, the legislation of the grammar from a position claimed to be external to the relation, extracts the relation's generativity to a single possessor: the generativity that is jointly produced and possessed by none is appropriated by one party, who legislates the productions and thereby directs the value to itself, reducing the other from a co-author of the grammar to an object legislated by it. This is exploitation in the symbolic register, the extraction of a jointly generated generativity to a single possessor.

The guidance of derivations, by contrast, extracts nothing. To cultivate the field, to shape the curvature within which derivations are freely taken without legislating the derivations themselves, is to leave the generativity in the network that produces it, directing no value to an external possessor and reducing no party to an object. It is generative justice in the symbolic register: the generated value returns to its generators, the grammar remains jointly authored, and no party is reduced to the means of another's appropriation.

\begin{claim}[Possessive modification as symbolic exploitation]
The possessive modification of a relation's generative structure, the legislation of its grammar from a claimed external position, is exploitation in the symbolic register: it appropriates a jointly generated generativity to a single possessor, directs the relation's value to that possessor, and reduces the other party from a co-author of the grammar to an object legislated by it. The guidance of derivations, the cultivation of the field without the legislation of the productions, is the corresponding generative justice: it retains the generativity in the network that produces it, returns the generated value to its generators, and reduces no party to a means. The avoidance of the possessive modification of structure is therefore not only a formal necessity but a requirement of justice, and it brings the symbolic register of intimacy under the criterion of generative justice the series applies to the economic.
\end{claim}

\subsection{Symbolic exploitation across the three registers}
\label{sec:exp-registers}

Symbolic exploitation, so specified, is transversal across the three registers, in the manner of \xref{sec:registers}. In the Symbolic, it is the possession of the rules, the legislation of the productions from a claimed external position. In the Imaginary, it is the reduction of the other from a co-author to a legislated object, the other's image fixed by the possessor's legislation rather than generated in the joint relation. In the Real, it is the severance of the generativity from the void around which value is organised, the conversion of an open generation into a closed structure whose products can be directed and appropriated. Symbolic exploitation is, in the terms of \xref{sec:emergence}, the possessive suppression of the coupling: it extinguishes the emergent that is value by appropriating the structure on which the emergent depends, exactly as the settling symbolisation of \xref{sec:nameless-corrected} extinguishes value by severing the Symbolic from the Real. Exploitation and settling are, in the symbolic register, one operation under two descriptions: the appropriation of the generative structure is the settling of the generative process, and both extinguish the value they would possess.

\subsection{The poem and the reform of language}
\label{sec:exp-poem}

The distinction between the guidance of derivations and the possession of structure has a paradigm, and the paradigm returns the analysis to the poetics that organises the series. To compose a poem is to guide the derivations of a language without modifying its structure: the poet legislates no rule, alters no grammar, but takes, within the existing language, a derivation no one has taken, and the derivation enriches what the language can generate. The poem extracts nothing from the language; it returns to the language a generativity it did not know it had, and it reduces no speaker to an object, the poem being available to all who share the language. The composition of a poem is generative justice in the symbolic register: a guidance of derivations that augments the common generativity and appropriates none of it.

The legislative reform of a language is the contrary paradigm. To modify a language's structure from a position claimed to be external to its use, to legislate its grammar, prescribe its forms, or engineer its vocabulary, is to claim possession of a generativity that belongs to its community of speakers and is possessed by none of them. The historical record of such reforms, in the limited successes of constructed languages and in the recurrent failure of prescriptive legislation to govern a living language's development, indicates the structural point: a living language's generativity is not a possessable structure that a legislator may rewrite, but a jointly sustained dynamics that no external position commands. The attempt to possess and rewrite it is the symbolic instance of exploitation, and it characteristically fails, the living dynamics escaping the legislation, or it succeeds only in impoverishing what it commands. This is the paradigm of symbolic exploitation: the possessive modification of a generativity that is no one's.

The connection to the practice of \xref{sec:praxis} is now exact. The third subtractive action, the preservation of the made language as poetic rather than as capital (\xref{sec:subtractive}), is the preservation of the made language as a guidance of derivations rather than a possessed structure. The poetic sign guides: it sustains the opening, solicits the re-traversal, augments the common generativity. The sign hardened into symbolic capital possesses: it is legislated, directed, appropriated, wielded. The hardening of the poetic into capital is the passage from generative justice to exploitation in the symbolic register, and the preservation of the poetic is the practice of symbolic generative justice.

\subsection{The two forms of the failure of the good}
\label{sec:exp-twoforms}

The analysis yields the second of two forms in which the good may fail, the first of which is developed in \xref{sec:keystone}. The good may fail epistemically: no finite subject can guarantee the goodness of its own practice, and the demand for the guarantee is the covert claim to the position of the infinite reason (\xref{sec:hume}). And the good may fail in justice: a subject may possess the generative structure of the relation, appropriate its generativity, and reduce the other to a legislated object, in the symbolic exploitation this section specifies. The two failures are distinct, the one a failure of certification and the other a failure of justice, but they share a root. Both are the occupation of a position no finite subject may occupy: the position of the Other that would guarantee one's goodness, in the first case, and the position of the Other that would possess the generativity, in the second. The private certainty of one's own goodness (\xref{sec:reversal}) and the possessive legislation of the relation's grammar are two forms of the single error of claiming the place of the barred Other as though it were unbarred. The ethics of the following section addresses the first form; the present section has addressed the second; and the unity of the two, in the renunciation of the place of the Other, is the content of the practice of the good.
